% Minimal retained preamble for lane-p2 cards taken from
% sources/logic-fl02-conservation-of-ramsey-s-theorem-f/pinned-original-source.tex
% (Conservation of Ramsey's theorem for pairs, arXiv:2402.11616v3).
%
% Copied verbatim from the author's own file:
%   line 2      : \documentclass{article}
%   lines 16-25 : the author's theorem declarations
%   lines 46-88 : the author's macro definitions for the reverse-mathematics principles
%   line 91     : \card
%   lines 104-107: the natural sums and products, which rely on mathabx (source line 103)
%   line 116    : \qt
%
% Omitted author preamble lines and why (nothing was renamed or redefined):
%   line 3 (graphicx), 7 (url), 9 (cancel): no figure, url or cancel use in the retained spans;
%   line 6 (cleveref): the assembler adds its own cleveref machinery, a second load clashes;
%   line 8 (\usepackage[utf8x]{inputenc}): inputenc with the utf8x option fights the forced
%       xelatex engine;
%   lines 28-37 (the rep@theorem machinery and \newreptheorem{maintheorem}): the retained spans
%       use no rep-theorem environment;
%   lines 40-43 (xcolor and the three author comment macros): not used by the retained spans;
%   lines 93-99 (the LS2 font encoding block for \smashtimes based on stix): the host TeX Live has
%       no stix directory and the author himself never uses the result -- line 105 maps \ntimes to
%       \dottimes, and line 110 shows \smashtimes already commented out in the source.
% mathabx is retained because \dotplus, \dottimes and \dotdiv come from it.
\documentclass{article}
\usepackage{amsmath,amssymb,amsthm}
\usepackage{mathabx}
\usepackage{cancel}
\usepackage{cleveref}

\newtheorem{theorem}{Theorem}
\numberwithin{theorem}{section}
\newtheorem{maintheorem}[theorem]{Main Theorem}
\newtheorem{lemma}[theorem]{Lemma}
\newtheorem{question}[theorem]{Question}
\newtheorem{proposition}[theorem]{Proposition}
\newtheorem{statement}[theorem]{Statement}
\newtheorem{remark}[theorem]{Remark}
\newtheorem{definition}[theorem]{Definition}
\newtheorem{corollary}[theorem]{Corollary}

\newcommand{\RCA}{\mathsf{RCA}}
\newcommand{\WKL}{\mathsf{WKL}}
\newcommand{\SRT}{\mathsf{SRT}}
\newcommand{\EM}{\mathsf{EM}}
\newcommand{\SEM}{\mathsf{SEM}}
\newcommand{\ADS}{\mathsf{ADS}}
\newcommand{\SCAC}{\mathsf{SCAC}}
\newcommand{\SADS}{\mathsf{SADS}}
\newcommand{\COH}{\mathsf{COH}}
\newcommand{\ACA}{\mathsf{ACA}}
\newcommand{\qvdash}{\operatorname{{?}{\vdash}}}
\newcommand{\nqvdash}{\operatorname{{?}{\nvdash}}}
\newcommand{\BSig}{\mathsf{B}\Sigma^0}
\newcommand{\ISig}{\mathsf{I}\Sigma^0}
\newcommand{\LSig}{\mathsf{L}\Sigma^0}
\newcommand{\BPi}{\mathsf{B}\Pi^0}
\newcommand{\IPi}{\mathsf{I}\Pi^0}
\newcommand{\LPi}{\mathsf{L}\Pi^0}
\newcommand{\IDel}{\mathsf{I}\Delta^0}
\newcommand{\BCDe}{\mathsf{BC}\Delta^0}
\newcommand{\WF}{\mathsf{WF}}
\newcommand{\RT}{\mathsf{RT}}
\newcommand{\BME}{\mathsf{BME}}
\newcommand{\TT}{\mathsf{TT}}
\renewcommand{\O}{\mathcal{O}}
\renewcommand{\L}{\mathcal{L}}
\newcommand{\C}{\mathcal{C}}
\newcommand{\R}{\mathcal{R}}
\renewcommand{\S}{\mathcal{S}}
\newcommand{\T}{\mathcal{T}}
\renewcommand{\P}{\mathcal{P}}
\newcommand{\M}{\mathcal{M}}
\newcommand{\Nc}{\mathcal{N}}
\renewcommand{\P}{\mathcal{P}}
\newcommand{\Q}{\mathcal{Q}}
\newcommand{\NN}{\mathbb{N}}
\newcommand{\PP}{\mathbb{P}}
\newcommand{\uh}[0]{{\upharpoonright}}
\newcommand{\card}{\operatorname{card}}
\newcommand{\nplus}{\dotplus}
\newcommand{\ntimes}{\dottimes}
\newcommand{\nminus}{\dotdiv}
\newcommand{\nsum}{\dot{\sum}}
\def\qt#1{``#1''}
\title{Every pair of Sigma\textasciicircum{}0\_1 formulas admits a maximal tree of split pairs satisfying it}
\author{Quentin Le Houérou, Ludovic Levy Patey, Keita Yokoyama}
\date{arxiv:2402.11616v3; CC BY 4.0}
\begin{document}
\maketitle
\noindent\textit{Editorial note.} The counted unit is the article's own Lemma lem:sads-wf-exists-maximal-tree at source lines 1194-1196 together with the proof environment that begins two lines after the statement ends (1198-1204); the pairing is the ordinary adjacent one.\par\medskip
\noindent\textit{Editorial note.} The retained context at source lines 1169-1183 is the author's own definition of a tree of split pairs, of its labelling and of a valid node; the retained context at source lines 1186-1192 is his own definition of satisfaction of a pair of Sigma\textasciicircum{}0\_1 formulas by a node and by a tree. Both definitions sit outside the surrounding proof environment of prop:sads-wf, which this lane does not count.\par\medskip
\noindent\textit{Editorial note.} The retained preamble reproduces verbatim the author's document class at source line 2, his theorem declarations at source lines 16-25, his macro definitions at source lines 46-88 and 91, and his natural-operation macros at source lines 104-107, and it retains mathabx and cleveref; the stix-dependent LS2 block and the inputenc load with the utf8x option are not retained.\par\medskip
\noindent\textit{Editorial note.} No retained line contains a citation, so the card carries no bibliography.\par\medskip
\noindent\textit{Editorial note.} The article's conservation theorems, its COH, WKL and SEM lemmas and its other SADS lemmas are not counted here.\par\medskip
\noindent\textit{Allowed standard prerequisites (editorial).}\begin{itemize}
\item Trees of split pairs and valid nodes. A tree attached to a condition is a binary tree whose nodes are labelled by split pairs extending the condition, with the left and right successors extending the two components respectively; a node is valid when along the left branches above it the ascending components lie inside U and along the right branches the descending components lie inside V; the author's own definition is retained at source lines 1169-1183.. Supplies the objects of the counted lemma, namely the tree, its labels and the maximality that the construction must enforce.
\item Satisfaction of a pair of Sigma\textasciicircum{}0\_1 formulas. A node labelled by a split pair satisfies a pair of Sigma\textasciicircum{}0\_1 formulas when the first formula holds of the ascending component at the rank of the node and the second holds of the descending component at the number of ones in the node, and a tree satisfies the pair when every node does; retained at source lines 1186-1192.. It is the property that the increasing sequence of trees must preserve at every stage, and it is what makes the limit tree satisfy the pair.
\end{itemize}\medskip
\begin{definition}
Fix a condition $(\sigma, \tau)$.

A \emph{$(\sigma, \tau)$-tree} is a binary tree $T$ labelled by a family of split pairs $(\sigma_{\rho}, \tau_{\rho})_{\rho \in T}$ such that for all $\rho \in T$:
\begin{itemize}
    \item $\sigma \preceq \sigma_{\rho}$ and $\tau \preceq \tau_{\rho}$ ;
    \item For every $\mu \in T$, if $\rho \cdot 0 \preceq \mu$ then $\sigma_{\rho} \preceq \sigma_{\mu}$ and if $\rho \cdot 1 \preceq \mu$ then $\tau_{\rho} \preceq \tau_{\mu}$ ;
\end{itemize}
%    For a node $\rho \in T$, we define $\pi(\rho)$ to be equal to $(\sigma_{\rho \uh i_0}, \tau_{\rho \uh j_0})$ where $i_0$ is the biggest index $i < |\rho|$ such that $\rho(i) = 0$ (if no such index exists then we take $\sigma$ instead of $\sigma_{\rho \uh i_0}$ in the definition of $\pi(\rho)$) and $j_0$ is the biggest index $j < |\rho|$ such that $\rho(j) = 1$ (simillarly, if no such index exists, take $\tau$ instead of $\tau_{\rho \uh j_0}$). 

%    A node is said to be \emph{valid} if $\pi(\rho)$ is a condition.
     %A node $\rho$ is said to be \emph{valid} if the biggest index $i$ such that $\rho(i) = 0$ satisfies $\sigma_{\rho \uh i} \subseteq U$ (if no such index exists this condition is verified) and if the biggest index $j$ such that $\rho(i) = 1$ satisfies $\tau_{\rho \uh i} \subseteq V$ (if no such index exists this condition is verified). Being valid is a $\Delta_2^0(Y)$ condition.

     A node $\rho$ is said to be \emph{valid} if for every $\mu \in T$ such that $\mu \cdot 0 \preceq \rho$, $\sigma_{\mu} \subseteq U$ and for every $\mu \in T$ such that $\mu \cdot 1 \preceq \rho$, $\tau_{\mu} \subseteq V$. Being valid is a $\Delta^0_2(Y)$ condition.
\end{definition}

\begin{definition}
    Let $T$ a $(\sigma,\tau)$-tree 
    \begin{itemize}
        \item A node $(\sigma_{\rho}, \tau_{\rho})$ in $T$ \emph{satisfies} a pair of $\Sigma^0_1(Y)$-formulas $(\varphi_0(G_0, x), \varphi_1(G_1,x))$ if $\varphi_0(\sigma, |\rho| - \card \rho)$ and $\varphi_1(\tau, \card \rho)$ holds (where $\card \rho$ is the cardinal of $\rho$ seen as a set).
        \item $T$ \emph{satisfies} $(\varphi_0(G_0, x), \varphi_1(G_1,x))$ if every node of $T$ satisfies it.
    \end{itemize}
\end{definition}

\begin{lemma}\label[lemma]{lem:sads-wf-exists-maximal-tree}
For every pair $(\varphi_0(G_0, x), \varphi_1(G_1,x))$ of $\Sigma^0_1(Y)$ formulas and every condition $(\sigma, \tau)$, there is a maximal $\Sigma^0_1$ $(\sigma, \tau)$-tree~$T$ satisfying it.
\end{lemma}

\begin{proof}
$T$ will be the limit of the following increasing sequence of $Y$-computable $(\sigma, \tau)$-trees $(T_s)_{s \in M}$ that all satisfy $(\varphi_0(G_0, x), \varphi_1(G_1,x))$ (and therefore $T$ will also be a $(\sigma, \tau)$-tree satisfying $(\varphi_0(G_0, x), \varphi_1(G_1,x))$). 

At stage $0$, let $T_0$ be the empty tree. It vacuously satisfies $(\varphi_0(G_0, x), \varphi_1(G_1,x))$.

At stage $s+1$, if there exists a split pair $(\sigma', \tau')$ with $|\sigma'|,|\tau'| \leq s$, a $\rho \in T_s$ and an $i < 2$ with $\rho \cdot i \notin T_s$ such that adding $\rho \cdot i$ to $T_s$ with the index $(\sigma', \tau')$ still yields a $(\sigma, \tau)$-tree satisfying $(\varphi_0(G_0, x), \varphi_1(G_1,x))$, then let $T_{s+1}$ be that tree (choose the smallest split pair possible to enforce the maximality of $T$). If no such element is found, simply go to the next stage.
\end{proof}

\end{document}
